% Options for packages loaded elsewhere
\PassOptionsToPackage{unicode}{hyperref}
\PassOptionsToPackage{hyphens}{url}
\PassOptionsToPackage{dvipsnames,svgnames,x11names}{xcolor}
%
\documentclass[
  11pt,
  a4paper,
]{article}
\usepackage{amsmath,amssymb}
\usepackage{iftex}
\ifPDFTeX
  \usepackage[T1]{fontenc}
  \usepackage[utf8]{inputenc}
  \usepackage{textcomp} % provide euro and other symbols
\else % if luatex or xetex
  \usepackage{unicode-math} % this also loads fontspec
  \defaultfontfeatures{Scale=MatchLowercase}
  \defaultfontfeatures[\rmfamily]{Ligatures=TeX,Scale=1}
\fi
\usepackage{lmodern}
\ifPDFTeX\else
  % xetex/luatex font selection
  \setmainfont[]{Latin Modern Roman}
  \setmathfont[]{Latin Modern Math}
\fi
% Use upquote if available, for straight quotes in verbatim environments
\IfFileExists{upquote.sty}{\usepackage{upquote}}{}
\IfFileExists{microtype.sty}{% use microtype if available
  \usepackage[]{microtype}
  \UseMicrotypeSet[protrusion]{basicmath} % disable protrusion for tt fonts
}{}
\makeatletter
\@ifundefined{KOMAClassName}{% if non-KOMA class
  \IfFileExists{parskip.sty}{%
    \usepackage{parskip}
  }{% else
    \setlength{\parindent}{0pt}
    \setlength{\parskip}{6pt plus 2pt minus 1pt}}
}{% if KOMA class
  \KOMAoptions{parskip=half}}
\makeatother
\usepackage{xcolor}
\usepackage[top=25mm,bottom=24mm,left=27mm,right=27mm]{geometry}
\setlength{\emergencystretch}{3em} % prevent overfull lines
\providecommand{\tightlist}{%
  \setlength{\itemsep}{0pt}\setlength{\parskip}{0pt}}
\setcounter{secnumdepth}{-\maxdimen} % remove section numbering
% Mathematics-focused preprint layout: monochrome, Latin Modern OpenType text and math.
% Latin Modern fonts are supplied by a standard TeX installation, not embedded in this source archive.
\usepackage{microtype}
\usepackage{setspace}
\usepackage{titlesec}
\usepackage{fancyhdr}
\usepackage{needspace}
\setstretch{1.13}
\setlength{\parindent}{1.15em}
\setlength{\parskip}{3.25pt plus 1.0pt minus .6pt}
\titlespacing*{\section}{0pt}{17pt plus 4pt minus 2pt}{7pt plus 2pt}
\titlespacing*{\subsection}{0pt}{13pt plus 3pt minus 1pt}{5pt plus 1pt}
\titlespacing*{\subsubsection}{0pt}{10pt plus 2pt minus 1pt}{3pt plus 1pt}
\titleformat{\section}{\normalfont\bfseries\large}{\thesection}{0.7em}{}
\titleformat{\subsection}{\normalfont\bfseries\normalsize}{\thesubsection}{0.6em}{}
\titleformat{\subsubsection}{\normalfont\bfseries\normalsize}{\thesubsubsection}{0.6em}{}
\AtBeginDocument{%
  \setlength{\abovedisplayskip}{12pt plus 3pt minus 2pt}%
  \setlength{\belowdisplayskip}{12pt plus 3pt minus 2pt}%
  \setlength{\abovedisplayshortskip}{9pt plus 2pt minus 1pt}%
  \setlength{\belowdisplayshortskip}{10pt plus 2pt minus 1pt}%
  \pagestyle{plain}%
}
\setlength{\emergencystretch}{4em}
\setcounter{secnumdepth}{0}
\ifLuaTeX
  \usepackage{selnolig}  % disable illegal ligatures
\fi
\usepackage{bookmark}
\IfFileExists{xurl.sty}{\usepackage{xurl}}{} % add URL line breaks if available
\urlstyle{same}

\usepackage{longtable,booktabs,array,calc}
\usepackage{fvextra}
\RecustomVerbatimEnvironment{verbatim}{Verbatim}{breaklines=true,breakanywhere=true,fontsize=\footnotesize}
\usepackage{upquote}

\hypersetup{
  pdftitle={Primitive norm-three CM dynamics on secp256k1},
  pdfauthor={Louis Desbiens (human initiator and curator)},
  pdfsubject={AI-developed mathematical research artifact; independently unverified candidate results},
  colorlinks=true,
  linkcolor={black},
  filecolor={black},
  citecolor={black},
  urlcolor={black},
  pdfcreator={LaTeX via pandoc}}

\title{Primitive norm-three CM dynamics on secp256k1}
\usepackage{etoolbox}
\makeatletter
\providecommand{\subtitle}[1]{%
  \apptocmd{\@title}{\par {\large #1 \par}}{}{}
}
\makeatother
\subtitle{A sixfold quotient orbit and the scope of algebraic scalar-map obstructions}
\author{}
\date{Public mathematical research artifact \textbar{} 31 October 2026}
\begin{document}
\maketitle

\textbf{Human initiator and curator:} Louis Desbiens\\
\textbf{AI mathematical research contributor:} OpenAI's ChatGPT\\
\textbf{Document classification:} Public AI-developed mathematical
research artifact -- independently unverified candidate results\\
\textbf{Institutional affiliation:} None claimed\\
\textbf{Manuscript license:} Creative Commons Attribution 4.0 International (CC BY 4.0)

\textbf{Status and reading guide.} This document records a
human-initiated, substantially AI-developed mathematical investigation.
The human curator directed the objectives but does not independently
certify its advanced mathematical arguments. The novel claims and
supplied proofs are candidate results awaiting independent specialist
verification; public availability does not certify correctness,
completeness or originality. This artifact investigates primitive
norm-three complex multiplication, an explicit quotient map and
restrictions on certain rational scalar maps. Specific secp256k1
arithmetic is accompanied by machine-rerunnable certificates; the
general rigidity application, conditional correspondence arguments and
any cryptographic interpretation need separate scrutiny. None of the
claims constitutes a general algorithmic lower bound or an ECDLP
speedup.

\section{Abstract}\label{abstract}

This manuscript examines the order-three automorphism of secp256k1 and a
degree-three endomorphism whose scalar eigenvalue is proposed to be
primitive modulo the prime group order. On the complete sixfold quotient
\(t=x^3\), the endomorphism induces the explicit rational map
\(F(t)=-(t+28)^3/(27t^2)\). The analysis includes its derivation, an
exact arithmetic order certificate, a proposed single-cycle description
of the quotient states, and the identity expressing \(F\) as a
functional square root of quotient tripling. This reframes the discrete
logarithm modulo sixfold symmetry as an orbit-index problem, without
giving a faster method to solve it. A proposed finite-set agreement
argument derives the restricted degree bound \((n-1)/72\) for rational
quotient maps implementing hidden scalar squaring. Separately, the cited
classical rigidity theorem for maps into abelian varieties constrains
globally rational point-output constructions on \(E^r\) to affine
hidden-scalar actions, subject to its hypotheses. The manuscript
explores commuting quotient maps and a conditional statement about
invariant correspondences; it claims neither a complete
positive-characteristic centralizer classification nor a computational
lower bound for arbitrary programs. The investigation combines explicit
secp256k1 dynamics with established geometric theory, while submitting
its new deductions for independent review.

\textbf{Keywords:} secp256k1; complex multiplication; Lattes map;
primitive orbit; elliptic discrete logarithm; rational map; abelian
variety.

\section{1. Setting and contribution}\label{setting-and-contribution}

A particularly simple endomorphism occurs within the usual \(j=0\)
structure of secp256k1. This article focuses on that object and the
limits of operations that remain rational maps of the generic curve. The
surrounding theory of efficient endomorphisms, Lattes maps, and
abelian-variety rigidity is classical {[}1--4{]}. The arithmetic order
certificate and explicit quotient formulas below make the specific
arithmetic claims independently checkable; they do not establish
historical priority for those formulas.

Let

\[
E/\mathbf F_p:\quad y^2=x^3+7,\qquad p=2^{256}-2^{32}-977.
\]

The prime order \(n\) is specified in hexadecimal by concatenating the
following two 32-digit words:

\begin{verbatim}
FFFFFFFFFFFFFFFFFFFFFFFFFFFFFFFE
BAAEDCE6AF48A03BBFD25E8CD0364141
\end{verbatim}

The standard base point is denoted \(G\). The domain parameters and
cofactor one are those of SEC 2 {[}1{]}. In particular,
\(E(\mathbf F_p)=\langle G\rangle\) and every nonidentity rational point
has order \(n\). Fix the conventional orientation of the cube roots
\(\beta\in\mathbf F_p\) and \(\lambda\in\mathbf F_n\) by concatenating
each pair of hexadecimal words below:

\begin{verbatim}
beta   = 7ae96a2b657c07106e64479eac3434e9
         9cf0497512f58995c1396c28719501ee
lambda = 5363ad4cc05c30e0a5261c028812645a
         122e22ea20816678df02967c1b23bd72
\end{verbatim}

Here leading zeroes are immaterial. They satisfy \(\beta^2+\beta+1=0\),
\(\lambda^2+\lambda+1=0\), and

\[
\phi(x,y)=(\beta x,y),\qquad \phi(G)=[\lambda]G.
\]

The identity on \(G\) is a finite group-arithmetic check using the SEC 2
coordinates and the ordinary chord-and-tangent law. Since \(G\)
generates the rational group, checking it once proves the scalar action
on every rational point. Its coordinates, for direct reproduction, are
the following hexadecimal pairs:

\begin{verbatim}
x(G) = 79be667ef9dcbbac55a06295ce870b070
       29bfcdb2dce28d959f2815b16f81798
y(G) = 483ada7726a3c4655da4fbfc0e1108a8
       fd17b448a68554199c47d08ffb10d4b8
\end{verbatim}

Define

\[
H=\{\pm1,\pm\lambda,\pm\lambda^2\},\qquad h=2+\lambda,\qquad M=(n-1)/6. \qquad (1.1)
\]

The letter \(h\) denotes a scalar in this article, not the cofactor
convention used in standards documents.

\section{2. A primitive endomorphism of degree
three}\label{a-primitive-endomorphism-of-degree-three}

\textbf{Proposition 2.1.} The endomorphism

\[
\Psi=[2]+\phi=1-\phi^2
\]

is separable of degree three, acts as multiplication by \(h\) on
\(E(\mathbf F_p)\), and satisfies

\[
\Psi^2=-[3]\phi^2. \qquad (2.1)
\]

\textbf{Proof.} The relation \(1+\phi+\phi^2=0\) gives the second
expression for \(\Psi\). In the Eisenstein order,
\(N(a+b\omega)=a^2-ab+b^2\), so \(N(2+\omega)=3\). Equivalently,
composing \(2+\phi\) with \(1-\phi\) gives \([3]\), and the explicit
formulas below show degree three. Since the characteristic is not three,
the map is separable. Its rational scalar action follows by evaluating
the endomorphism on the generator. Finally,
\((2+\phi)^2=3+3\phi=-3\phi^2\). \(\square\)

\textbf{Proposition 2.2 (primitive scalar).} For the fixed orientation
in Section 1,

\[
\operatorname{ord}_n(h)=n-1. \qquad (2.2)
\]

\textbf{Certificate.} The complete factorization is

\[
n-1=2^6\cdot3\cdot149\cdot631\cdot q_1q_2q_3. \qquad (2.3)
\]

Here the three larger prime factors are

\[
\begin{gathered}
q_1=107361793816595537,\qquad q_2=174723607534414371449,\\
q_3=341948486974166000522343609283189.
\end{gathered}
\]

The factors multiply exactly to \(n-1\). Appendix A prints a complete
recursive Lucas certificate, including the primality of these factors
and of \(n\). Its last row uses \(h\) as the witness and checks

\[
h^{n-1}\equiv1\pmod n,\qquad
\gcd(h^{(n-1)/q}-1,n)=1
\]

for each distinct prime factor \(q\) in (2.3). The criterion proved in
the appendix implies (2.2). All certificate data are included in this
article. \(\square\)

The certificate is a short exact computation, not an enumeration of
\(n-1\) points. Multiplication by a primitive element acts transitively
on \(\mathbf F_n^*\); consequently \(\Psi\) is one cycle on the
nonidentity rational points.

Changing from \(\lambda\) to its conjugate is an orientation change, not
a license to reuse every residue identity unchanged. The quotient map
below is insensitive to multiplication of the endomorphism by a sixfold
unit, but character normalizations in later work are not.

\section{3. The complete sixfold
quotient}\label{the-complete-sixfold-quotient}

\textbf{Proposition 3.1.} On nonidentity rational points,
\(t(P)=x(P)^3\) is a complete invariant of the \(H\)-orbit. Every such
orbit has six elements, and there are \(M\) quotient states.

\textbf{Proof.} Negation fixes \(x\), and \(\phi\) multiplies \(x\) by
\(\beta\), so \(t\) is invariant. If \(x(Q)^3=x(P)^3\), then
\(x(Q)=\beta^j x(P)\) for some \(j\). The curve equation gives
\(y(Q)=\pm y(P)\), hence \(Q=\pm\phi^jP\). No nonidentity point of the
prime-order group is fixed by a nontrivial element of \(H\): the
fixed-point conditions would force two- or three-torsion. Thus every
orbit has six elements. \(\square\)

\textbf{Proposed Theorem 3.2 (quotient dynamics).} The map induced by
\(\Psi\) is

\[
x(\Psi P)=-\frac{\beta}{3}\frac{x(P)^3+28}{x(P)^2},\qquad
F(t)=-\frac{(t+28)^3}{27t^2}. \qquad (3.1)
\]

On the set of \(M\) nonidentity rational quotient states, \(F\) is a
single cycle. As rational maps on \(\mathbf P^1\),

\[
F^{\circ2}=L_3, \qquad (3.2)
\]

where \(L_3\) is the map induced by multiplication by three on \(E/H\).

\textbf{Proof.} Write \(\Psi P=P-\phi^2P\). The joining slope is
\(-2y/[x(\beta^2-1)]\). Using \((\beta^2-1)^2=-3\beta^2\) and
\(1+\beta^2=-\beta\) in the chord formula gives the first identity in
(3.1). Cubing it gives the second. The rational functions extend across
their affine poles as morphisms of the projective line.

The quotient identifies scalar multiplication by the unique subgroup of
order six of \(\mathbf F_n^*\). Since \(h\) is primitive, its class
generates \(\mathbf F_n^*/H\), a cyclic group of order \(M\). This
proves the cycle statement. Equation (2.1) differs from multiplication
by three by a unit of \(H\), which disappears on the quotient; this
proves (3.2) as a generic rational identity, not merely a finite-set
equality. \(\square\)

Neither \(x=0\) nor an affine two-torsion point belongs to the
nonidentity rational subgroup. Therefore the displayed affine
denominators do not create exceptional states within this cycle.

\section{4. Orbit index and hidden
scalar}\label{orbit-index-and-hidden-scalar}

Let \(Q=[d]G\) with \(d\neq0\). There is an exponent \(\kappa\) modulo
\(n-1\) with \(d=h^\kappa\). Passing to the quotient yields

\[
t(Q)=F^\kappa(t(G)),\qquad \kappa\pmod M. \qquad (4.1)
\]

\textbf{Proposed Corollary 4.1 (exact reformulation).} Solving the
quotient orbit-index problem in (4.1) determines the six candidates

\[
d_j=h^{\kappa+jM}\pmod n,\qquad 0\leq j<6. \qquad (4.2)
\]

Exactly one satisfies \([d_j]G=Q\). Conversely, knowledge of \(d\)
determines the target point and its quotient orbit position, although
computing \(\log_h d\) is itself a multiplicative discrete logarithm.

The useful computational equivalence is in the direction from a
recovered quotient shift to the six checked scalar candidates. One
should not assert an efficient reduction from an arbitrary scalar-output
ECDLP solver to the integer orbit index without accounting for the
additional finite-field logarithm \(\log_h d\). The phrase ``hidden
shift'' refers here to the orbit relation; (4.1) and (4.2) specify its
exact meaning.

A collision between the two walks \(F^A(t(G))\) and \(F^B(t(Q))\) gives
\(\kappa=A-B\pmod M\). Such a relation is useful only if \(A\) and \(B\)
are known. A compact formula for a low-norm endomorphism does not
automatically provide its exponent relative to \(h\).

For example, a jump induced by \(a+b\phi\) has degree \(a^2-ab+b^2\) and
scalar \(a+b\lambda\). Its primitive-orbit increment is

\[
\log_h(a+b\lambda)\pmod M,
\]

not the easily computed pair \((a,b)\). This is a known-increment
bottleneck: the map and its orbit displacement are different inputs. No
sub-square-root orbit-index algorithm follows from the existence of
\(F\).

\section{5. Critical portrait and rational
predecessors}\label{critical-portrait-and-rational-predecessors}

The logarithmic derivative of \(F\) is

\[
\frac{F'(t)}{F(t)}=\frac{t-56}{t(t+28)}.
\]

Its ramification points are \(-28\) with local degree three, \(0\) with
local degree two, and \(56\) with local degree two. Their forward
portrait is

\[
-28\longmapsto0\longmapsto\infty\longmapsto\infty,
\qquad 56\longmapsto-7\longmapsto-7. \qquad (5.1)
\]

The normalized coordinate \(u=-t/7\) gives

\[
\mathcal G(u)=\frac{(4-u)^3}{27u^2}. \qquad (5.2)
\]

The quotient has ramification signature \((2,3,6)\), as expected for
this elliptic quotient. This places the map in the classical Lattes
framework {[}4{]}; it does not by itself prove every possible
positive-characteristic classification assertion.

For a finite target \(v\), a predecessor solves

\[
t^3+(84+27v)t^2+2352t+21952=0.
\]

The discriminant is

\[
\Delta(v)=-(9072v)^2\,21(v+7). \qquad (5.3)
\]

For a nonidentity quotient state, \(v+7\) is a nonzero square, while
\(-21\) is a nonsquare in the secp256k1 field; the latter follows from
\(\chi_2(-3)=1\) and \(\chi_2(7)=-1\). Indeed \(p\equiv1\pmod3\)
supplies a square root of \(-3\), while a square root of \(7\) would
give a rational point \((0,\sqrt7)\) of order three, contrary to the
prime group order. The cubic is unramified there. Since the quotient
cycle supplies a rational predecessor, its remaining factor is
irreducible quadratic: a nonsquare discriminant excludes three rational
roots. Thus one rational branch is singled out over \(\mathbf F_p\), and
the other two are conjugate over its quadratic extension. This
arithmetic fact does not make the orbit index easy to recover; the
predecessor is already the inverse of a permutation on an enormous
cycle.

\section{6. A degree obstruction for hidden
squaring}\label{a-degree-obstruction-for-hidden-squaring}

Define \(t_d=t([d]G)\) for \(d\neq0\). The rule \(t_d\mapsto t_{d^2}\)
is well defined because multiplication of \(d\) by a unit of \(H\)
changes \(d^2\) by another unit of \(H\).

\textbf{Proposed Theorem 6.1 (finite-set degree lower bound).} If a
nonconstant rational function \(R\) is defined on every nonidentity
quotient state and satisfies

\[
R(t_d)=t_{d^2}\qquad(d\in\mathbf F_n^*),
\]

then

\[
\deg R\geq\frac{M}{12}=\frac{n-1}{72}. \qquad (6.1)
\]

\textbf{Proof.} At each of the \(M\) states,

\[
R\circ F(t_d)=t_{h^2d^2}=F^{\circ2}\circ R(t_d).
\]

Writing \(D=\deg R\), the two rational maps have degrees \(3D\) and
\(9D\). They cannot be identical, because their positive degrees differ.
Two distinct rational maps of degrees \(a,b\) on \(\mathbf P^1\) have at
most \(a+b\) distinct agreement points: clear their homogeneous
denominators and count the zeros of the resulting nonzero homogeneous
polynomial. Thus \(M\leq12D\). \(\square\)

A constant \(R\) cannot satisfy the proposed rule: the squaring map on
the cyclic quotient has more than one value. The theorem therefore
covers all rational candidates. It is an exact finite-field statement,
but it concerns \textbf{formal degree}, not evaluation cost. For
example, \(F^{\circ k}\) has degree \(3^k\) and a composition circuit of
depth \(k\). The generic point-output restriction below is a separate
argument and does not use a degree-to-time inference.

\section{7. Classical rigidity and point-output
constructions}\label{classical-rigidity-and-point-output-constructions}

\textbf{Theorem 7.1 (classical).} Every rational map from an abelian
variety to an abelian variety extends to a morphism and is a
homomorphism followed by a translation. In particular,

\[
\mathcal F:E^r\dashrightarrow E\quad\Longrightarrow\quad
\mathcal F(P_1,\ldots,P_r)=T+\sum_{j=1}^r\alpha_j(P_j), \qquad (7.1)
\]

where \(\alpha_j\in\operatorname{End}(E)\) and \(T\in E\).

\textbf{Proof.} The extension theorem for maps from a smooth variety to
an abelian variety {[}3, Theorem 3.1{]} extends the rational map to a
morphism. Subtract its value \(T\) at the origin to obtain \(g\) with
\(g(0)=0\). The morphism \(D(P,Q)=g(P+Q)-g(P)-g(Q)\) vanishes on both
axes. The rigidity lemma {[}3, Theorem 3.4{]} makes it identically zero,
so \(g\) is a homomorphism. Restricting \(g\) to each factor gives the
endomorphisms \(\alpha_j\), and adding the coordinate inclusions gives
(7.1). If the original map is defined over the base field, uniqueness of
extension and of the restrictions preserves that field of definition.
\(\square\)

\textbf{Proposed Corollary 7.2 (hidden-scalar action).} If
\(\mathcal F\) is defined over \(\mathbf F_p\), its action on secp256k1
rational points is

\[
(d_1,\ldots,d_r)\longmapsto c+\sum_{j=1}^rs_jd_j\pmod n. \qquad (7.2)
\]

Indeed, \(T=[c]G\) and \(\alpha_j(G)=[s_j]G\). Therefore no such
globally rational point-output map implements \([d]G\mapsto[d^2]G\), or
\(([d_1]G,[d_2]G)\mapsto[d_1d_2]G\), on the whole prime-order group.
These nonlinear functions cannot coincide with an affine function on
that group.

The application includes any arithmetic circuit whose output coordinates
satisfy the elliptic-curve equation as an identity in the generic
function field of \(E^r\). Its formal degree and compactness do not
evade the theorem. Intermediate rational maps through auxiliary
varieties also do not evade it when the final composition is genuinely a
rational map \(E^r\dashrightarrow E\).

This is a structural statement about one computational class. It is not
a generic-group lower bound, and not a proof that a faster algorithm for
ECDLP cannot exist. The generic-group theorem of Shoup {[}5{]} concerns
a different model and must not be substituted for a proof about a
particular coordinate representation.

\section{8. What the rigidity statement does not
say}\label{what-the-rigidity-statement-does-not-say}

\subsection{8.1. Finite-field identities need not define generic curve
maps}\label{finite-field-identities-need-not-define-generic-curve-maps}

For a fixed finite field, exponentiation and even a multiplicative
character can be expressed by polynomial coordinate functions. Thus
``uses a character'' is not literally the same as ``is not algebraic.''
The decisive issue in Theorem 7.1 is whether the \textbf{point-valued
output} defines a rational map on the generic curve, not whether the
implementation uses multiplication or exponentiation.

For a rational affine point \(Q=(a,b)\), the polynomial

\[
\delta_Q(x,y)=\bigl(1-(x-a)^{p-1}\bigr)\bigl(1-(y-b)^{p-1}\bigr)
\]

is its indicator on \(\mathbf F_p^2\). Linear combinations of these
indicators can interpolate arbitrary finite tables of affine output
coordinates. Such outputs can lie on \(E\) at all relevant rational
inputs without satisfying the curve equation generically. This
observation illustrates the scope limitation; it is not an efficient
interpolation algorithm for an unknown discrete logarithm. Projective
output coordinates can handle tables containing \(O\).

Likewise, selecting roots, testing bits, comparing representatives, or
branching may lead outside the generic point-output class. The theorem
does not bound the information or computation obtainable from such
operations.

\subsection{8.2. Affine leaves are not running-time
bounds}\label{affine-leaves-are-not-running-time-bounds}

Suppose a branch of a program is a globally rational point-output map
and agrees with hidden squaring on its assigned scalar inputs. Its
scalar expression is \(sd+c\), so it can agree with \(d^2\) at at most
two field elements. Covering all nonzero scalars therefore requires at
least \((n-1)/2\) such affine pieces. If equality is required only
modulo \(H\), each branch gives at most twelve roots by considering
\(d^2=u(sd+c)\) for the six units \(u\), hence at least \((n-1)/12\)
pieces are needed.

These are bounds on the \textbf{number of distinct affine pieces}, under
the stated branch model. They are not exponential time or circuit lower
bounds. A program making \(k\) binary decisions can have \(2^k\) paths,
and shared subcomputations may further change its representation size.
Neither leaf counting nor generic-map rigidity resolves arbitrary
finite-field branching complexity.

\subsection{8.3. Squaring and
multiplication}\label{squaring-and-multiplication}

An exact hidden-squaring operation would imply hidden multiplication by
polarization:

\[
[d_1d_2]G=[2^{-1}]\bigl(S(P+Q)-S(P)-S(Q)\bigr),
\]

where \(P=[d_1]G\), \(Q=[d_2]G\), and \(S([d]G)=[d^2]G\). The inverse of
two is taken modulo \(n\). This identifies a useful research target, but
does not assert that such an operation has been constructed.

\section{9. Pairings and the field-of-definition
obstruction}\label{pairings-and-the-field-of-definition-obstruction}

The arithmetic certificate in Appendix A also gives

\[
\operatorname{ord}_n(p)=M. \qquad (9.1)
\]

Appendix A verifies both \(p^M=1\pmod n\) and \(p^{M/q}\neq1\pmod n\)
for every distinct prime divisor \(q\) of \(M\). The field containing a
nontrivial \(n\)th root of unity must therefore have extension degree
divisible by \(M\).

The Weil pairing on two points of the same cyclic rational subgroup is
trivial by alternation and bilinearity. CM endomorphisms defined over
the base field keep that subgroup inside itself. Moving the outputs to
an auxiliary low-dimensional variety does not change the elementary
condition \(n\mid p^m-1\) for a nontrivial \(\mu_n\)-valued pairing over
\(\mathbf F_{p^m}\). This is a field-of-definition obstruction for such
pairings. It does not classify every possible use of auxiliary varieties
or every non-pairing reduction.

\section{10. Commuting quotient maps and invariant
correspondences}\label{commuting-quotient-maps-and-invariant-correspondences}

\textbf{Proposition 10.1.} Every nonzero Eisenstein endomorphism
\(\alpha=a+b\phi\) induces a rational map \(R_\alpha\) on \(E/H\), and
\(R_\alpha\circ F=F\circ R_\alpha\). Multiplication of \(\alpha\) by a
sixfold unit does not change \(R_\alpha\).

\textbf{Proof.} The function \(t\circ\alpha\) is invariant under \(H\),
since \(\alpha\) commutes with \(\phi\) and negation. The invariant
function field is \(\mathbf F_p(t)\), so
\(t\circ\alpha=R_\alpha\circ t\). Commutation of \(\alpha\) with
\(\Psi\) gives \(R_\alpha F t=F R_\alpha t\). Since \(t\) is dominant,
the rational maps agree. Multiplication by a unit changes no quotient
value. \(\square\)

This proves an inclusion in the rational centralizer. It is not a proof
that the listed family exhausts that centralizer. A complete converse in
positive characteristic requires, among other things, treatment of
inseparability and of the lift from the quotient; the complex-analytic
classification in {[}4{]} is not used as a substitute for that argument.

\textbf{Proposition 10.2 (a conditional correspondence statement).} Let
\(Z\subset E\times E\) be an irreducible projective curve with a
nonconstant projection to \(E\). Suppose a self-map of \(Z\) induces on
its smooth normalization \(\widetilde Z\) a separable self-map of degree
\(D>1\). Then \(\widetilde Z\) has genus one, and over an algebraic
closure the image of \(\widetilde Z\) in \(E\times E\) is a translate of
an elliptic subgroup.

\textbf{Proof.} Riemann--Hurwitz gives \(2g-2\geq D(2g-2)\), so
\(g\leq1\). Genus zero is impossible: every map from \(\mathbf P^1\) to
an abelian variety is constant {[}3, Corollary 3.8{]}. Choose a point on
the genus-one normalization as origin. Its map to \(E\times E\),
translated to send this origin to zero, is a homomorphism by {[}3,
Corollary 3.6{]}. Its image is a one-dimensional abelian subvariety.
Undoing the translation proves the claim. \(\square\)

For a component preserved only by an iterate, the hypothesis must be
checked for that iterate. The translation need not vanish. This
proposition does not classify every correspondence on the quotient line,
nor does it convert a multivalued correspondence into a point-valued
algorithm.

\section{11. Interpretation and
conclusion}\label{interpretation-and-conclusion}

The exploratory investigation used numerical evaluations to test the
initial intuition behind the quotient dynamics. The cycle assertion
itself follows from the complete arithmetic certificate, not from
observing a long orbit. Appendix A also supplies the checks for the
embedding degree and the orientation identities.

The positive object is an explicit primitive degree-three quotient
dynamics. The principal restriction is the classical affine form of
globally rational point-output maps. Their combination narrows a class
of proposed constructions without proving a general security statement.
Large formal degree, finite-field interpolation, branching, and discrete
labels must be analyzed in their own computational models. No method
beating the generic square-root scale is established here.

\section{Appendix A. Complete arithmetic
certificate}\label{appendix-a.-complete-arithmetic-certificate}

This appendix gives all data needed to certify the scalar order; it does
not rely on an external computation file. The prime coordinate modulus
and the curve parameters are the published SEC 2 parameters. The order
\(n\), the three large factors of \(n-1\), and every auxiliary prime
used below are certified by the following elementary criterion.

\textbf{Lucas criterion.} Suppose the complete factorization
\(q-1=\prod t^{e_t}\) is known, every \(t\) has already been proved
prime, and an integer \(a\) satisfies

\[
a^{q-1}\equiv1\pmod q,\qquad
\gcd(a^{(q-1)/t}-1,q)=1\quad\text{for every }t\mid q-1.
\]

Then \(q\) is prime and \(a\) has order \(q-1\) modulo \(q\). Indeed,
for every prime divisor \(\ell\) of \(q\), the displayed conditions
force the order of \(a\) modulo \(\ell\) to contain the full power of
every prime dividing \(q-1\). Thus \(q-1\mid\ell-1\), so \(\ell\geq q\)
and \(\ell=q\).

Use the parameters \(n,\lambda\) of Section 1 and \(h=\lambda+2\). The
labels in the table mean

\[
q_1=107361793816595537,\qquad q_2=174723607534414371449,
\] \[
q_3=341948486974166000522343609283189,
\] \[
r_0=29047611873442575647497758179.
\]

Starting with the prime \(2\), the rows are in increasing numerical
order. Each row states a complete factorization and a Lucas witness.
Integer multiplication, modular exponentiation, and the Euclidean
algorithm verify every row. No probabilistic primality test is required.

\begin{longtable}[]{@{}
  >{\raggedright\arraybackslash}p{(\columnwidth - 4\tabcolsep) * \real{0.3000}}
  >{\raggedright\arraybackslash}p{(\columnwidth - 4\tabcolsep) * \real{0.3000}}
  >{\centering\arraybackslash}p{(\columnwidth - 4\tabcolsep) * \real{0.4000}}@{}}
\toprule\noalign{}
\begin{minipage}[b]{\linewidth}\raggedright
\(q\)
\end{minipage} & \begin{minipage}[b]{\linewidth}\raggedright
Factorization of \(q-1\)
\end{minipage} & \begin{minipage}[b]{\linewidth}\centering
Witness \(a\)
\end{minipage} \\
\midrule\noalign{}
\endhead
\bottomrule\noalign{}
\endlastfoot
\(3\) & \(2\) & \(2\) \\
\(5\) & \(2^{2}\) & \(2\) \\
\(7\) & \(2\cdot 3\) & \(3\) \\
\(11\) & \(2\cdot 5\) & \(2\) \\
\(13\) & \(2^{2}\cdot 3\) & \(2\) \\
\(17\) & \(2^{4}\) & \(3\) \\
\(19\) & \(2\cdot 3^{2}\) & \(2\) \\
\(23\) & \(2\cdot 11\) & \(5\) \\
\(29\) & \(2^{2}\cdot 7\) & \(2\) \\
\(37\) & \(2^{2}\cdot 3^{2}\) & \(2\) \\
\(41\) & \(2^{3}\cdot 5\) & \(6\) \\
\(59\) & \(2\cdot 29\) & \(2\) \\
\(67\) & \(2\cdot 3\cdot 11\) & \(2\) \\
\(73\) & \(2^{3}\cdot 3^{2}\) & \(5\) \\
\(97\) & \(2^{5}\cdot 3\) & \(5\) \\
\(109\) & \(2^{2}\cdot 3^{3}\) & \(6\) \\
\(113\) & \(2^{4}\cdot 7\) & \(3\) \\
\(149\) & \(2^{2}\cdot 37\) & \(2\) \\
\(199\) & \(2\cdot 3^{2}\cdot 11\) & \(3\) \\
\(293\) & \(2^{2}\cdot 73\) & \(2\) \\
\(461\) & \(2^{2}\cdot 5\cdot 23\) & \(2\) \\
\(631\) & \(2\cdot 3^{2}\cdot 5\cdot 7\) & \(3\) \\
\(797\) & \(2^{2}\cdot 199\) & \(2\) \\
\(1409\) & \(2^{7}\cdot 11\) & \(3\) \\
\(1871\) & \(2\cdot 5\cdot 11\cdot 17\) & \(14\) \\
\(2011\) & \(2\cdot 3\cdot 5\cdot 67\) & \(3\) \\
\(2731\) & \(2\cdot 3\cdot 5\cdot 7\cdot 13\) & \(3\) \\
\(2861\) & \(2^{2}\cdot 5\cdot 11\cdot 13\) & \(2\) \\
\(4051\) & \(2\cdot 3^{4}\cdot 5^{2}\) & \(10\) \\
\(9349\) & \(2^{2}\cdot 3\cdot 19\cdot 41\) & \(2\) \\
\(16699\) & \(2\cdot 3\cdot 11^{2}\cdot 23\) & \(3\) \\
\(28181\) & \(2^{2}\cdot 5\cdot 1409\) & \(2\) \\
\(85831\) & \(2\cdot 3\cdot 5\cdot 2861\) & \(3\) \\
\(120233\) & \(2^{3}\cdot 7\cdot 19\cdot 113\) & \(3\) \\
\(305873\) & \(2^{4}\cdot 7\cdot 2731\) & \(3\) \\
\(1627771\) & \(2\cdot 3\cdot 5\cdot 29\cdot 1871\) & \(3\) \\
\(4681609\) & \(2^{3}\cdot 3\cdot 97\cdot 2011\) & \(23\) \\
\(44706919\) & \(2\cdot 3\cdot 797\cdot 9349\) & \(6\) \\
\(545358713\) & \(2^{3}\cdot 41\cdot 59\cdot 28181\) & \(5\) \\
\(297159362677\) & \(2^{2}\cdot 3^{2}\cdot 11\cdot 461\cdot 1627771\) &
\(2\) \\
\(q_1\) & \(2^{4}\cdot 16699\cdot 85831\cdot 4681609\) & \(3\) \\
\(q_2\) & \(2^{3}\cdot 17\cdot 59\cdot 4051\cdot 120233\cdot 44706919\)
& \(3\) \\
\(r_0\) & \(2\cdot 293\cdot 305873\cdot 545358713\cdot 297159362677\) &
\(2\) \\
\(q_3\) & \(2^{2}\cdot 3^{3}\cdot 109\cdot r_0\) & \(2\) \\
\(n\) & \(2^{6}\cdot 3\cdot 149\cdot 631\cdot q_1\cdot q_2\cdot q_3\) &
\(h\) \\
\end{longtable}

The last row uses \(h\) itself. It simultaneously proves that \(n\) is
prime and that \(h\) has order \(n-1\). Its seven proper-order tests are
precisely those for the distinct factors \(2,3,149,631,q_1,q_2,q_3\); in
particular, the certificate is not a traversal of the curve group.

The following additional checks use the same stated parameters. They can
be reproduced by the short exact-arithmetic code printed below;
exponentiation is repeated squaring. Here hexadecimal strings are
concatenated before conversion.

\begin{verbatim}
from math import gcd
p = 2**256 - 2**32 - 977
n = int('FFFFFFFFFFFFFFFFFFFFFFFFFFFFFFFE'
        'BAAEDCE6AF48A03BBFD25E8CD0364141', 16)
beta = int('7ae96a2b657c07106e64479eac3434e9'
           '9cf0497512f58995c1396c28719501ee', 16)
lam = int('5363ad4cc05c30e0a5261c028812645a'
          '122e22ea20816678df02967c1b23bd72', 16)
h, M = lam + 2, (n - 1) // 6
factors = [2, 3, 149, 631, 107361793816595537,
           174723607534414371449,
           341948486974166000522343609283189]
assert (beta*beta + beta + 1) % p == 0
assert (lam*lam + lam + 1) % n == 0
assert pow(h, n-1, n) == 1
for t in factors:
    assert gcd(pow(h, (n-1)//t, n)-1, n) == 1
assert pow(p, M, n) == 1
for t in factors:
    if M % t == 0:
        assert pow(p, M//t, n) != 1
assert pow(h, M, n) == (-lam*lam) % n
assert M % 3 == 1 and p % 9 == 7
assert pow((beta*pow(81, -1, p)) % p, (p-1)//3, p) == 1
assert pow(19+beta, (p-1)//3, p) == beta
assert pow((-beta*pow(3, -1, p)) % p, (p-1)//3, p) == 1
assert (n-2)**2 > (9*(3**78-1))**2*p
assert (n-2)**2 <= (9*(3**79-1))**2*p
assert (n-3**39)**2 > (6*3**39*(3**39-1))**2*p
assert (n-3**40)**2 <= (6*3**40*(3**40-1))**2*p
\end{verbatim}

The code uses ordinary arbitrary-precision integers; negative modular
exponents denote inverses. The four squared inequalities are used only
with positive sides. They therefore certify, without floating-point
rounding, the thresholds in the character-sum applications. These finite
arithmetic certificates support the explicitly stated parameters and
identities; they are not a substitute for the geometric proofs.

\section{Research provenance, contributions, and verification
status}\label{research-provenance-contributions-and-verification-status}

\subsection{Origin and development of the
investigation}\label{origin-and-development-of-the-investigation}

This mathematical investigation was initiated and directed by Louis
Desbiens and developed through extensive interaction with OpenAI's
ChatGPT. Louis Desbiens identified the original research questions, set
the objectives, guided the investigation, assessed its progress at a
conceptual level, and selected the material for public dissemination.

ChatGPT contributed substantially to mathematical exploration, including
alternative formulations, derivation of algebraic identities,
formulation and examination of conjectures and proposed proofs,
identification of potential counterexamples and exceptional cases,
design and execution of symbolic and numerical checks, organization of
relevant literature, and drafting and revision of the text. These
contributions went well beyond linguistic or editorial assistance.

\subsection{Attribution and scientific
authorship}\label{attribution-and-scientific-authorship}

Louis Desbiens is identified as the \textbf{human initiator and
curator}, not as an independent mathematical authority certifying the
results. OpenAI's ChatGPT is identified as the \textbf{AI mathematical
research contributor}, not as an accountable conventional scholarly
coauthor. These designations describe the provenance without implying
that the human curator personally derived or independently verified all
the arguments.

\subsection{Mathematical verification
status}\label{mathematical-verification-status}

The candidate mathematical contributions have not undergone independent
specialist validation, formal proof verification, or an exhaustive
originality and priority review. In particular, independent examination
is needed of the complete arithmetic certificate and fixed endomorphism
orientation; the quotient-cycle deduction; hypotheses and use of
classical rigidity; scope of the finite-set degree obstruction and
conditional correspondences. Locally displayed proofs are proposed
arguments. Established results cited from public literature should be
assessed separately from new deductions; neither an AI-generated proof
nor a successful computational check is by itself an independently
certified theorem. Individual statements may be correct, partly correct,
incomplete, or incorrect.

\subsection{Responsibility and purpose of public
release}\label{responsibility-and-purpose-of-public-release}

Louis Desbiens assumes responsibility for the decision to release this
collection and for a good-faith account of its origin, contributors and
verification status. He does \textbf{not} claim to have personally
established or independently certified the correctness, completeness or
originality of the proposed mathematical results. The material is made
available for independent examination, criticism, reproducibility and
possible further development; publication does not constitute scientific
endorsement. This disclosure does not eliminate ordinary
responsibilities associated with distributing public material.

\subsection{Institutional independence, funding, and competing
interests}\label{institutional-independence-funding-and-competing-interests}

No institutional affiliation or endorsement is claimed. The
investigation was undertaken independently, without institutional
funding or resources. No external research funding was received. The
human curator reports no competing interests.

\subsection{Computational evidence and source
availability}\label{computational-evidence-and-source-availability}

Numerical evaluations helped develop the initial intuition. The public
repository contains the available scripts, archived results, and a
claim-by-claim verification coverage guide. These are scoped finite
arithmetic, symbolic, and numerical checks, not independent verification
of the general mathematical results. Historical exploratory calculations
for which independent source data are unavailable are not treated as
reproducible evidence.

\section{References}\label{references}

{[}1{]} Certicom Research, \emph{SEC 2: Recommended Elliptic Curve
Domain Parameters}, version 2.0, 27 January 2010, Section 2.4.1.
\url{https://www.secg.org/sec2-v2.pdf}.

{[}2{]} R. P. Gallant, R. J. Lambert, and S. A. Vanstone, ``Faster point
multiplication on elliptic curves with efficient endomorphisms,''
\emph{CRYPTO 2001}, LNCS 2139, 190--200, 2001. DOI:
10.1007/3-540-44647-8\_11. Public conference paper:
\url{https://www.iacr.org/archive/crypto2001/21390189.pdf}. For an
openly accessible treatment of endomorphisms and their arithmetic use,
see {[}6{]}.

{[}3{]} J. S. Milne, ``Abelian varieties,'' in \emph{Arithmetic
Geometry}, Springer, 1986, pp.~103--150. Corrected public author
version, 2 January 2022, Section 3 (Theorem 3.1, Theorem 3.4, Corollary
3.6): \url{https://www.jmilne.org/math/xnotes/AVs.pdf}.

{[}4{]} J. Milnor, \emph{On Lattes Maps}, arXiv:math/0402147v2, 2004.
\url{https://arxiv.org/abs/math/0402147v2}. Used for context on complex
elliptic-quotient dynamics, not a positive-characteristic
classification.

{[}5{]} V. Shoup, ``Lower bounds for discrete logarithms and related
problems,'' \emph{EUROCRYPT 1997}, LNCS 1233, 256--266, 1997. Public
author version: \url{https://www.shoup.net/papers/dlbounds1.pdf}. Its
generic-group model is distinct from the restricted geometric model
considered here.

{[}6{]} S. D. Galbraith, \emph{Mathematics of Public Key Cryptography},
Cambridge University Press, 2012. Public author version and errata:
\url{https://www.math.auckland.ac.nz/}\textasciitilde sgal018/crypto-book/crypto-book.html.
\end{document}
